\begin{titlepage}
  \centering
  \vspace*{30mm}
  {\Huge\bfseries MARS laboratory companion\par}
  \vspace{8mm}
  {\Large Experiments and independent assembly programs\par}
  \vspace{18mm}
  {\Large\MARSAuthor\par}
  \vfill
  {\MARSEdition\par}
\end{titlepage}

\chapter*{Using the laboratory companion}
\markboth{Using the laboratory companion}{Using the laboratory companion}
\addcontentsline{toc}{chapter}{Using the laboratory companion}
\label{revision:using-lab-book}

Use each numbered lab with the corresponding chapter of \emph{Programming the
  machine}. The main book develops the concept and worked reasoning. This companion
explains the experiment, the state to inspect, the failure to investigate and the
related program to write independently. The additional companion assignments are
placed after the concepts they need. They are opportunities to apply those ideas,
not hidden prerequisites for the later codec work.

Each lab names its public source files and checks. Keep a copy of your own
implementation as it develops. A complete worked experiment shows one behavior
under stated conditions; an independent starter leaves the related problem for
you to solve. Use the stated interface to decide whether a different
implementation is correct. Register allocation, helper labels and a reference
instruction sequence are not substitutes for that interface.

For an interactive lab, choose \textbf{Tools, Import MARS Lab...}. Open the
student ZIP or use the course's HTTPS archive link, choose \textbf{Preview lab},
review the edition, then choose \textbf{Import lab}. Its questions and actions
appear in Learn. Answer a prediction in that panel; the source pointer identifies
the instruction under discussion. For a Debug observation, use its named stops
and the debugger's actual selected frame to identify the paused instruction.
The lab explains which values to inspect and what they mean.

Use \textbf{Continue Downloaded MARS Lab} to return to the selected edition.
\textbf{Choose Downloaded MARS Lab Activity...} lets you return to guided work or
open an independent exercise when course policy permits it. Both routes keep the
same source and checking history. A hint is a way to investigate a difficulty;
you can first attempt the independent task without opening it. Instructors can
control prerequisites, hints, access and attempts. Private instructor checks,
when used for a course grade, remain separate from the public feedback here.

A successful execution observation belongs to the source revision that produced
it. After changing that source or a required earlier activity, renew the affected
checks. Changing a program back later does not make an old execution current.
Read a failure message as a counterexample to your reasoning, make a purposeful
change, and test it again. Explain why the change addresses that counterexample.

\begin{samepage}
  The imported work is below
  \path{mars-labs/packages/<id>/<version>/<sha256>} in the project. The assignment
  names paths relative to its own package root, the folder containing
  \texttt{package.json}. Book-checkout commands instead use the public repository
  root. Keep those two working locations distinct when following a command. The
  specified Run working directory also determines the meaning of relative filenames
  used by a MIPS program.\par
\end{samepage}

The complete source files also support ordinary Run and Debug outside a guided
sequence. Keep the supplied caller and storage when an assignment's checks use
them, then replace the designated unfinished function or program. Follow the
chapter's preparation instructions for device or file experiments. Preserve
outputs before comparing a changed experiment, and use a new report path where
the checker refuses to replace an existing report.

The glossary provides concise definitions. The subject index leads to the labs
that exercise each topic. The public files include deterministic input generators
and public checks; instructor references, alternatives and private cases belong
to the separate instructor distribution.
